% I. Предметна област и обосновка на избора.
% Обединява предишните раздели „литературен преглед“ и „анализ на алтернативите“.
\section{Предметна област и обосновка на избора}
\label{sec:literature}

\subsection{Теоретична рамка}

Работното определение за \term{облачни изчисления} описва достъп при поискване до споделен резерв
от ресурси и извежда пет характеристики: самообслужване, широк мрежов достъп, обединяване на
ресурсите, бърза еластичност и измерена услуга~\cite{nist800145}. Последните две са пряко
относими тук: еластичността е това, което проверява синтетичното натоварване, а измерената услуга
прави сравнението по цена възможно изобщо. Еластичността струва повече от средната утилизация,
защото премахва нуждата от оразмеряване по върховото натоварване~\cite{armbrust2010view}. Трите
модела на обслужване са IaaS (възли, мрежа и съхранение), PaaS (доставчикът поема средата за
изпълнение) и SaaS (готово приложение). При управляваните услуги за контейнери
границата между IaaS и PaaS е размита, и точно това прави сравнението нетривиално: еднакво
наименована услуга може да лежи от различни страни на границата.

Технологичната основа на обединяването на ресурси е \term{виртуализацията}. Хардуерната
виртуализация чрез хипервайзор изолира пълни виртуални машини със собствено
ядро~\cite{barham2003xen}, докато контейнеризацията, тоест виртуализация на ниво операционна
система, споделя едно ядро между изолирани пространства от имена и стартира многократно по-бързо
за сметка на по-слаба изолация~\cite{merkel2014docker}. Изборът между двете определя и времето за
реакция при мащабиране, и профила на риска.

Управлението на много еднакви възли води до \term{клъстерните технологии}: планировчик, който
държи $N$ реплики живи, каквато е и цялата нужда на този проект от
оркестрация~\cite{verma2015borg}. Предшественик са \term{GRID технологиите} за обединяване на
ресурси между отделни организации~\cite{foster2001grid}; разликата с облака е в собствеността и
таксуването, не в техническата задача. Обектните хранилища жертват част от семантиката на
файловата система срещу висока достъпност, като двама доставчици правят различен избор по оста
достъпност спрямо съгласуваност~\cite{decandia2007dynamo,calder2011azurestorage}. Сигурността се
управлява по модела на споделената отговорност, чието следствие тук е отказът от дълготрайни
ключове в полза на \term{управлявана идентичност}: възелът получава краткотраен маркер от самата
платформа и никакъв секрет не се записва в хранилището с изходния текст.

\subsection{Направени избори}

Изборите са четири и всеки е оценен по три критерия, формулирани преди избора: наличие на
съпоставим еквивалент при двамата доставчици, възможност описанието да остане едно, и пригодност
на измерванията за сравнение. За изпълнение на работника са разгледани виртуални машини (пълен
контрол, но стартиране от порядъка на минути), управляван Kubernetes и управлявани услуги за
контейнери. За описание на инфраструктурата отпадат императивните скриптове и собствените средства
на доставчиците, които работят само при своя доставчик; избран е Terraform, съответно OpenTofu,
защото при него разликата се локализира в отделни модули. Три от изборите бяха преразгледани при
реализацията, и Табл.~\ref{tab:choices} показва и първоначалния, и окончателния вариант
\tabref{tab:choices}: премълчаването на първия би направило обосновката съчинена след факта.

\begin{table}[H]
\caption{Направени избори, първоначални и след реализацията}
\label{tab:choices}
\centering
\small
\begin{tabular}{@{}L{3.0cm}L{2.8cm}L{2.9cm}L{5.0cm}@{}}
\toprule
\textbf{Решение} & \textbf{Първоначално} & \textbf{Окончателно} & \textbf{Причина и цена на промяната} \\
\midrule
Изпълнение на работника & Управляван Kubernetes, EKS и AKS &
ECS с Fargate и Container Apps &
Услугата не използва нищо от Kubernetes освен $N$ живи реплики на едно име, а вторият слой на мащабиране замъглява измерваната величина \\
\addlinespace
Описание на инфраструктурата & Terraform или OpenTofu & Без промяна & Няма \\
\addlinespace
Генератор на натоварване & k6 & Собствен, на C++ &
Нула задължителни външни зависимости. Цена: около 160 реда. Печалба: избор между затворен и отворен контур \\
\addlinespace
Събиране на метрики & Prometheus и Grafana &
Форматът на Prometheus, със собствено събиране &
Изборът на \emph{формат} запазва същественото: всеки скрейпър чете метриките без промяна в приложението. Цена: няма готово табло \\
\bottomrule
\end{tabular}
\end{table}

Общото между трите промени е критерий, който не беше формулиран ясно в началото: изграждане с нула
външни зависимости.
